\documentclass[conference]{IEEEtran}
\IEEEoverridecommandlockouts

\usepackage{cite}
\usepackage{amsmath,amssymb,amsfonts}
\usepackage{algorithm}
\usepackage{algpseudocode}
\usepackage{graphicx}
\usepackage{booktabs}
\usepackage{array}
\usepackage{textcomp}
\usepackage{xcolor}
\usepackage[hidelinks]{hyperref}
\usepackage{url}

\renewcommand{\arraystretch}{1.12}

\begin{document}

\title{AI-Driven Multi-Objective Orchestration and\\
Adaptive Scheduling for Noise-Aware\\
Distributed Quantum Computing}

\author{\IEEEauthorblockN{Author Name}
\IEEEauthorblockA{\textit{Department / Institution} \\
\textit{Affiliation}\\
City, Country \\
email@domain}
}

\maketitle

\begin{abstract}
Distributed quantum computing (DQC) couples several noisy intermediate-scale
quantum (NISQ) processors so that circuits exceeding the capacity or fidelity of
any single device can be executed cooperatively. Operating such a fleet requires
an orchestration layer that decides which processor(s) run which circuit, in what
order, subject to objectives that are simultaneously in tension: delivered
fidelity, decoherence-bounded deadlines, throughput, and load balance. Classical
schedulers optimize none of the quantum-specific objectives, and a naive
fidelity-greedy policy ignores deadlines and collapses under contention. We
present AMOS, an AI-driven, multi-objective orchestrator with an adaptive,
noise-aware scheduling core. AMOS (i)~formulates placement as the maximization of
a weighted multi-objective host score that \emph{satisfices} fidelity---rewarding
tasks that clear a fidelity threshold rather than maximizing fidelity
unconditionally---while jointly optimizing deadline feasibility, speed, and load;
(ii)~orders tasks by least laxity to protect coherence budgets under overload;
and (iii)~adapts online to calibration drift through a lightweight reliability
estimator that derates processors whose realized fidelity falls below their
datasheet. We evaluate AMOS in a discrete-event simulator of eight heterogeneous
QPUs with decoherence, entanglement-mediated circuit distribution, and
calibration drift. AMOS matches the best baseline when the fleet is uncontended
and improves goodput by up to $\sim\!2\times$ under heavy overload while
sustaining a $>\!99\%$ deadline-hit rate. A weight-sensitivity study shows the
delivered-fidelity/goodput operating point is robust across the full
satisfice-to-maximize range (variation $<\!1$\,pp), which justifies the
satisficing default. An ablation shows the structured multi-objective policy,
not the learner, drives the gain in stationary conditions; online adaptation
earns its place only under non-stationarity, where it recovers a consistent
fraction of the goodput lost after a device-degradation event. We report these
findings---including the limits of the learning component---transparently and
release the simulator for reproduction.
\end{abstract}

\begin{IEEEkeywords}
distributed quantum computing, quantum orchestration, multi-objective scheduling,
NISQ, decoherence, entanglement distribution, online adaptation, noise-aware
scheduling
\end{IEEEkeywords}

% ===================================================================
\section{Introduction}
% ===================================================================
Individual quantum processors are limited in qubit count and, more
fundamentally, in the fidelity and coherence time available to a single
device~\cite{barral2025review}. Distributed quantum computing (DQC) mitigates
these limits by interconnecting multiple processors---within a chassis, a data
center, or an eventual quantum network---so that a large circuit can be
partitioned and executed across nodes using entanglement and classical
communication~\cite{caleffi2024survey,chehimi2022qin}; distributed execution has
recently been demonstrated across an optical network
link~\cite{main2025optical}. As quantum cloud
providers move toward multi-tenant operation over pools of heterogeneous
devices~\cite{nguyen2024drlq}, they need an \emph{orchestration layer}: a
control plane that admits a stream of quantum tasks (circuits) and maps them onto
a heterogeneous QPU fleet over time.

Orchestrating a quantum fleet is not classical job scheduling with new labels.
Four properties make off-the-shelf policies inadequate:

\begin{itemize}
\item \textbf{Fidelity is a first-class outcome.} Without error correction, a
NISQ execution returns a result whose usefulness is governed by process fidelity;
two placements that both ``complete'' can differ from useful to worthless.
\item \textbf{Deadlines are physical.} A circuit must finish within a
coherence-limited budget; gate error and decoherence jointly cap the useful depth
on a given device.
\item \textbf{Distribution is not free.} Executing across $k$ nodes introduces
non-local gates realized by entanglement (gate teleportation), each consuming an
EPR pair and adding latency and a multiplicative fidelity
penalty~\cite{piveteau2024knitting,du2024cutting}.
\item \textbf{Devices are non-stationary.} Calibration drifts; a processor that
looks best on its datasheet can be transiently or persistently
unreliable~\cite{nguyen2025qfor}.
\end{itemize}

These properties place the orchestrator in a genuinely \emph{multi-objective}
regime---delivered fidelity, deadline satisfaction, throughput, and load balance
must be traded against one another---and in a \emph{non-stationary} one, where the
datasheet the orchestrator relies on may be stale. Classical
throughput/load-oriented schedulers ignore fidelity, while a fidelity-greedy
policy ignores deadlines and contention and, as we show, collapses under load.

\noindent\textbf{Contributions.}
\begin{enumerate}
\item We formulate noise-aware DQC orchestration over heterogeneous NISQ
processors as multi-objective placement plus deadline-aware ordering, with a
\emph{goodput} objective coupling a decoherence-bounded deadline and a fidelity
threshold (Section~\ref{sec:model}).
\item We propose \textbf{AMOS}, an AI-driven orchestrator whose scheduling core
combines least-laxity ordering, a \emph{satisficing} multi-objective host score,
and an online reliability estimator that adapts to calibration drift
(Section~\ref{sec:method}).
\item We build a discrete-event simulator modeling decoherence,
entanglement-mediated distribution, and per-device drift, and evaluate AMOS
against four baselines with a controlled ablation
(Sections~\ref{sec:setup}--\ref{sec:results}).
\item We report results honestly: AMOS improves goodput by up to $\sim\!2\times$
under overload; its operating point is robust to the fidelity-weighting knob; and
online learning is negligible under stationarity but consistently helpful under
non-stationarity. We discuss the model's limitations explicitly
(Section~\ref{sec:limitations}).
\end{enumerate}

% ===================================================================
\section{Related Work}
% ===================================================================
\textbf{DQC foundations, surveys, and hardware.} The DQC landscape is mapped by
recent surveys spanning single-QPU limits to high-performance quantum
computing~\cite{caleffi2024survey,barral2025review}, reviews of recent
progress~\cite{knorzer2025review}, and taxonomies of the path toward
fault-tolerant DQC~\cite{larasati2025ftdqc}. The networked substrate is framed by
the quantum-internet vision and its building
blocks~\cite{chehimi2022qin,meddeb2025qinternet}. Crucially, DQC is no longer only
theoretical: recent hardware demonstrations couple processors with real-time
classical communication~\cite{carrera2024realtime} and even execute a distributed
algorithm across an optical network link~\cite{main2025optical}, making the
control-plane questions we study concrete.

\textbf{Circuit cutting, knitting, and multi-core mapping.} Distribution is made
concrete by circuit cutting/knitting, which decomposes a large circuit into
subcircuits at a classical/entanglement overhead that grows with the number of
cuts~\cite{piveteau2024knitting,du2024cutting}. Subsequent work tightens this
overhead for wire cuts~\cite{harada2024wire}, multiple two-qubit
unitaries~\cite{schmitt2025cutting}, and clustered simulation~\cite{harrow2025optimal}.
On the compiler side, modular compilation frameworks~\cite{ferrari2023modular},
qubit-assignment heuristics for multi-core
architectures~\cite{escofet2023hungarian}, and learned circuit
partitioning~\cite{pastor2024partitioning} target the \emph{intra-job} mapping.
We adopt an abstracted form of the resulting distribution cost in our execution
model and treat cutting/mapping as orthogonal to the \emph{across-job, over-time}
orchestration problem.

\textbf{Quantum cloud scheduling and orchestration.} Closest to our work, several
recent efforts schedule quantum jobs across a fleet. Bahrani et
al.~\cite{bahrani2024resource} study resource management and circuit scheduling
for DQC interconnect networks, Chandra et al.~\cite{chandra2025netops} schedule
the underlying network operations, Ferrari et al.~\cite{ferrari2024execution}
address execution management of DQC jobs, and Du et al.~\cite{du2024cutting}
co-optimize circuit cutting and scheduling with dynamic EPR pairs. Learning-based
orchestrators are increasingly common: DRLQ places tasks with a deep Q-network to
cut completion time~\cite{nguyen2024drlq}, QFOR adds fidelity awareness to a deep
reinforcement-learning orchestrator~\cite{nguyen2025qfor}, and Luo et al.\ apply
policy-gradient RL to adaptive job scheduling in quantum clouds~\cite{luo2025rl}.
Reviews and simulation toolkits situate these
efforts~\cite{nguyen2024review,nguyen2024iquantum}.

\textbf{Positioning.} AMOS differs from the above in two respects. First, its
scheduling core is a deliberately lightweight, interpretable multi-objective
\emph{satisficing} policy with an online reliability derate---no deep network, no
reward shaping---so its behavior is auditable and its weights map directly onto
operator objectives. Second, we report the value of the learned component
\emph{honestly}: an ablation shows it is negligible in stationary conditions and
helps only, and modestly, under non-stationarity---a distinction the deep-RL
orchestrators above do not isolate. To our knowledge, we are the first to couple a
decoherence-bounded deadline and a fidelity threshold into a single goodput
objective for across-fleet orchestration and to quantify when adaptation actually
pays.

% ===================================================================
\section{System Model and Problem Formulation}
\label{sec:model}
% ===================================================================

\subsection{Processors}
The fleet is a set $\mathcal{Q}=\{q_1,\dots,q_M\}$ of heterogeneous QPUs. Each
$q_j$ has qubit capacity $n_j$, nominal per-gate error $\varepsilon_j$, per-layer
time $\tau_j$, coherence time $T_j$ (a $T_2$-like budget), and a
calibration-instability factor $\delta_j$. The datasheet exposes
$(n_j,\varepsilon_j,\tau_j,T_j)$ to the orchestrator; $\delta_j$ and
moment-to-moment fluctuations are \emph{not} observable a priori.

\subsection{Tasks}
Tasks arrive as a stream. Task $i$ has arrival $a_i$, width $w_i$, depth $L_i$,
gate count $g_i$, and absolute deadline $d_i$. Arrivals are Poisson with rate set
by an offered-load parameter $\rho$ relative to aggregate fleet capability, so
larger $\rho$ means more contention. Deadlines are $d_i=a_i+\sigma_i\hat{t}_i$ for
a nominal estimate $\hat{t}_i$ and laxity $\sigma_i\in[2,6]$, modeling
coherence-bounded urgency.

\subsection{Execution and fidelity}
For a monolithic placement ($w_i\le n_j$) on $q_j$, execution time is
$t_{ij}=L_i\tau_j$ and the estimated process fidelity is
\begin{equation}
F_{ij} = \underbrace{(1-\varepsilon_j)^{g_i}}_{\text{gate errors}} \cdot
\underbrace{\exp\!\left(-t_{ij}/T_j\right)}_{\text{decoherence}} .
\label{eq:mono}
\end{equation}
When $w_i$ exceeds every free single device, the task is distributed over a set
$S$ of $k=|S|$ processors, introducing $c_i=\lceil\beta L_i(k-1)\rceil$ non-local
gates; each consumes an EPR pair of link fidelity $F_{\text{link}}$ and adds
latency $t_{\text{epr}}$:
\begin{align}
t^{\text{dist}}_i &= L_i \max_{j\in S}\tau_j + c_i t_{\text{epr}},\\
F^{\text{dist}}_i &= \Big(\textstyle\prod_{j\in S}(1-\varepsilon_j)^{g_i/k}\Big)
e^{-t^{\text{dist}}_i/\min_{j\in S}T_j} F_{\text{link}}^{\,c_i}.
\label{eq:dist}
\end{align}
Equations~\eqref{eq:mono}--\eqref{eq:dist} are intentionally first-order models
(Section~\ref{sec:limitations}); they capture the qualitative couplings---error
accumulation with gate count, decoherence with wall-clock time, and a compounding
entanglement penalty with the degree of distribution.

\subsection{Calibration drift}
The \emph{realized} outcome differs from the datasheet. At execution,
$\varepsilon_j$ and $1/T_j$ are each perturbed by an independent lognormal factor
with log-standard-deviation $\sigma_{\text{drift}}\delta_j$, so unstable devices
($\delta_j\!>\!1$) fluctuate more. A degradation event may persistently inflate a
device's error rate mid-stream (Section~\ref{sec:setup}) \emph{without} updating
the datasheet the orchestrator sees.

\subsection{Multi-objective goodput objective}
For each dispatched task we define objective components at placement: a fidelity
margin $m_i=F_i-F_{\text{th}}$, a deadline slack $\ell_i=d_i-f_i$ (with realized
finish $f_i$), execution speed, and incremental fleet load. A task \emph{succeeds}
iff it clears both the deadline and the fidelity threshold:
\begin{equation}
\text{success}_i = \mathbb{1}\!\left[f_i \le d_i\right]\cdot
\mathbb{1}\!\left[F_i \ge F_{\text{th}}\right].
\end{equation}
The orchestrator maximizes \emph{goodput}
$G=\frac{1}{N}\sum_{i}\text{success}_i$, the scalar that couples the fidelity and
deadline objectives; delivered fidelity, deadline-hit rate, and utilization are
reported as the constituent objectives. Even the offline, static-fidelity version
of this assignment is NP-hard (it generalizes deadline-constrained
unrelated-machine scheduling), motivating an online heuristic.

% ===================================================================
\section{AMOS: Adaptive Multi-Objective Orchestrator}
\label{sec:method}
% ===================================================================
AMOS runs at every orchestration event (arrival or completion). It maps a subset
of pending tasks onto free processors through three mechanisms
(Algorithm~\ref{alg:amos}).

\subsection{Least-laxity ordering}
Pending tasks are considered in increasing laxity
$\ell_i=d_i-t_{\text{now}}-\min_j t_{ij}$, an earliest-deadline-first discipline
that protects tasks at risk of exhausting their coherence budget. This directly
serves the deadline objective and is the main reason AMOS resists overload where
fidelity-greedy (arrival-ordered) policies do not.

\subsection{Satisficing multi-objective placement}
For a task and candidate host-set $h$ (a fitting single device, or one
distributed set), AMOS scores
\begin{equation}
s(i,h)=w_F\,\Phi\!\big(\hat{F}_{ih}\,r_{ih}\big)+w_D\,\psi_{ih}
+w_T\frac{t^{\min}_i}{t_{ih}}-w_L\,\bar{u}_h,
\label{eq:score}
\end{equation}
where $\hat{F}_{ih}$ is datasheet fidelity from
\eqref{eq:mono}--\eqref{eq:dist}, $r_{ih}$ is a learned reliability derate
(Section~\ref{sec:learn}), $\psi_{ih}\!\in\!\{-1,+1\}$ flags deadline
feasibility, $t^{\min}_i$ is the fastest feasible time, and $\bar{u}_h$ is an
exponential-moving load estimate. The fidelity mapping $\Phi$ blends a
satisficing and a maximizing view via an orchestration knob
$\lambda\!\in\![0,1]$:
\begin{equation}
\Phi(F)=(1-\lambda)\,\underbrace{\big(1+e^{-\gamma(F-F_{\text{th}})}\big)^{-1}}_{\text{satisfice}}
+\lambda\,\underbrace{F}_{\text{maximize}} .
\label{eq:phi}
\end{equation}
At $\lambda\!=\!0$ (default) the fidelity term saturates once the threshold is
safely cleared, so AMOS does not spend a scarce clean-but-slow device on a task
that is already safe elsewhere; at $\lambda\!=\!1$ it maximizes delivered fidelity
like a fidelity-greedy policy. A host-set whose best option cannot meet the
deadline is skipped rather than wasted. The weight vector
$(w_F,w_D,w_T,w_L)$ exposes the multi-objective trade-off to the operator; we use
$(1.6,1.3,0.7,0.25)$ throughout and study sensitivity in
Section~\ref{sec:results}.

\subsection{Online reliability adaptation}
\label{sec:learn}
Datasheet fidelity $\hat{F}$ is systematically optimistic for unstable devices.
AMOS maintains, per discretized task-state $x_i$ (width/depth buckets) and
processor $j$, two exponential moving averages: realized success
$\mathrm{rel}(x,j)$ and datasheet-predicted success $\mathrm{exp}(x,j)$. The
derate is their smoothed ratio
\begin{equation}
r_{ij}=\min\!\Big(1.1,\ \tfrac{\mathrm{rel}(x_i,j)+\eta}{\mathrm{exp}(x_i,j)+\eta}\Big),
\end{equation}
driving $<\!1$ for devices that underperform their spec and $=\!1$ for unseen
pairs (default trust). This is a lightweight, interpretable alternative to a deep
policy; it needs no reward shaping and adapts within tens of executions. Setting
$r_{ij}\!\equiv\!1$ gives the ablation \textbf{AMOS-S}.

\begin{algorithm}[t]
\caption{AMOS: one orchestration event}
\label{alg:amos}
\begin{algorithmic}[1]
\State \textbf{Input:} time $t$, pending $\mathcal{P}$, free hosts $\mathcal{F}$
\State sort $\mathcal{P}$ by ascending laxity $\ell_i$
\For{each task $i \in \mathcal{P}$}
  \State $\mathcal{C}\gets$ fitting single hosts in $\mathcal{F}$, else one
         distributed set
  \If{$\mathcal{C}=\varnothing$} \textbf{continue} \EndIf
  \For{each $h \in \mathcal{C}$}
     \State score $s(i,h)$ via \eqref{eq:score}--\eqref{eq:phi} with derate
     $r_{ih}$
  \EndFor
  \State $h^\star\gets\arg\max_h s(i,h)$
  \If{$t+t_{ih^\star}>d_i$} \textbf{continue} \Comment{skip hopeless}\EndIf
  \State assign $i\!\to\!h^\star$; mark busy; update load EMA
\EndFor
\State \textbf{on completion of} $i$ \textbf{on} $h$:
       update $\mathrm{rel},\mathrm{exp}$ for $(x_i,h)$
\end{algorithmic}
\end{algorithm}

Each event costs $O(|\mathcal{P}|\,|\mathcal{F}|)$ score evaluations, negligible
relative to circuit execution for a modest fleet.

% ===================================================================
\section{Experimental Setup}
\label{sec:setup}
% ===================================================================
\textbf{Fleet.} We model $M{=}8$ heterogeneous QPUs spanning the NISQ design
space: capacities $12$--$65$ qubits, per-gate errors $6\!\times\!10^{-4}$ to
$9\!\times\!10^{-3}$, layer times $3.0$--$8.0\!\times\!10^{-7}$\,s, and $T_2$ from
$7.0\!\times\!10^{-5}$ to $3.0\!\times\!10^{-4}$\,s. Two nominally attractive
devices are made calibration-\emph{unstable} ($\delta_j{=}2.4$--$2.6$) and one
clean device is made \emph{stable} ($\delta_j{=}0.6$), so datasheet ranking and
realized reliability diverge. Fixed physics: $F_{\text{th}}{=}0.50$,
$F_{\text{link}}{=}0.94$, $t_{\text{epr}}{=}2\,\mu$s, $\beta{=}0.12$,
$\sigma_{\text{drift}}{=}0.35$.

\textbf{Workload.} Widths are mostly $\le\!12$ with medium/large tails;
$\sim\!12\%$ of tasks exceed the largest free single device and require
distribution. Depths are $6$--$55$ layers, $g_i\!\approx\!0.5w_iL_i$. The
feasibility ceiling (fraction runnable above threshold on the best device) is
$\approx\!0.70$, and placement is decisive: $\approx\!69\%$ feasible on the
cleaner half of the fleet versus $\approx\!47\%$ on the noisier half. This is a
deliberately non-degenerate regime in which orchestration can matter
(Section~\ref{sec:limitations}).

\textbf{Baselines.} \emph{Random}; \emph{Round-Robin}; \emph{Least-Latency}
(FCFS, fastest free device); \emph{Max-Fidelity} (FCFS, highest datasheet
fidelity). We also report \textbf{AMOS-S} (no learning).

\textbf{Protocol.} Each configuration uses $N{=}1500$ tasks and $15$ seeds; all
policies see the \emph{same} workload and \emph{same} drift realization per seed
(paired comparison). We sweep offered load $\rho\in\{0.4,\dots,2.0\}$ and report
mean and s.d.\ of goodput, deadline-hit rate, mean delivered fidelity, and
utilization. The full simulator and orchestrator code accompany the paper; all
runs are seeded, so every figure and table reproduces exactly. Concretely,
\texttt{final\_run.py} generates Figs.~\ref{fig:goodput},~\ref{fig:bars},
and~\ref{fig:nonstat} with Tables~\ref{tab:metrics} and~\ref{tab:nonstat};
\texttt{pareto\_run.py} with \texttt{make\_sens\_fig.py} generate
Fig.~\ref{fig:sens} and Table~\ref{tab:sens}.

% ===================================================================
\section{Results}
\label{sec:results}
% ===================================================================

\subsection{Goodput vs.\ offered load}
Figure~\ref{fig:goodput} and Table~\ref{tab:metrics} summarize the sweep. When
uncontended ($\rho{=}0.4$) deadlines rarely bind and fidelity-greedy placement is
already near the feasibility ceiling: Max-Fidelity and AMOS-S tie at $0.63$
goodput, AMOS within noise. AMOS earns its advantage as contention rises. At
$\rho{=}1.5$, AMOS-S reaches $0.491$ versus $0.431$ for Max-Fidelity ($+13.9\%$
relative, $+6.0$\,pp) and $0.450$ for Random. At heavy overload ($\rho{=}2.0$) the
separation is stark: Max-Fidelity, Least-Latency, and Round-Robin all collapse to
$\approx\!0.22$--$0.23$ as they miss deadlines and thrash slow devices, whereas
AMOS sustains $0.450$---roughly a $2\times$ goodput improvement over
fidelity-greedy and $+28\%$ over the best baseline.

\begin{figure}[t]
\centering
\includegraphics[width=\columnwidth]{fig_goodput.pdf}
\caption{Goodput vs.\ offered load $\rho$ (mean $\pm$ s.d., 15 seeds). AMOS and
its static ablation track the best baseline when uncontended and degrade
gracefully under overload, where fidelity- and latency-greedy policies collapse.}
\label{fig:goodput}
\end{figure}

\subsection{Where the gain comes from}
Table~\ref{tab:metrics} isolates the mechanism. Under load, AMOS holds a
deadline-hit rate of $\approx\!0.995$--$0.996$ while baselines fall to
$0.937$--$0.967$ at $\rho{=}1.5$: least-laxity ordering does most of the work.
Notably, AMOS's mean delivered fidelity is slightly \emph{lower} than
Max-Fidelity's at $\rho{=}1.0$ ($0.468$ vs.\ $0.478$) even though its goodput is
higher---the satisficing trade-off made explicit: AMOS accepts adequate
(above-threshold) fidelity to meet more deadlines, and goodput, not raw fidelity,
is the objective. Figure~\ref{fig:bars} shows the cross-objective picture at
$\rho{=}1.5$.

\begin{table}[t]
\caption{Objective components at $\rho{=}1.0$ and $\rho{=}1.5$ (mean, 15 seeds).
G=goodput, DL=deadline-hit, F=mean fidelity, U=utilization.}
\label{tab:metrics}
\centering
\small
\begin{tabular}{@{}l cccc@{}}
\toprule
\multicolumn{5}{c}{\textbf{Offered load $\rho=1.0$}}\\
\midrule
Policy & G & DL & F & U\\
\midrule
Random        & 0.480 & 0.986 & 0.459 & 0.824\\
Round-Robin   & 0.485 & 0.986 & 0.460 & 0.826\\
Least-Latency & 0.451 & 0.981 & 0.440 & 0.840\\
Max-Fidelity  & 0.511 & 0.988 & \textbf{0.478} & 0.826\\
AMOS-S        & \textbf{0.522} & \textbf{0.996} & 0.470 & 0.823\\
AMOS          & 0.518 & \textbf{0.996} & 0.468 & 0.822\\
\midrule
\multicolumn{5}{c}{\textbf{Offered load $\rho=1.5$}}\\
\midrule
Policy & G & DL & F & U\\
\midrule
Random        & 0.450 & 0.967 & \textbf{0.475} & 0.982\\
Round-Robin   & 0.428 & 0.943 & 0.467 & 0.991\\
Least-Latency & 0.415 & 0.937 & 0.466 & 0.994\\
Max-Fidelity  & 0.431 & 0.941 & 0.469 & \textbf{0.996}\\
AMOS-S        & \textbf{0.491} & \textbf{0.995} & 0.474 & 0.991\\
AMOS          & 0.489 & \textbf{0.995} & 0.473 & 0.986\\
\bottomrule
\end{tabular}
\end{table}

\begin{figure}[t]
\centering
\includegraphics[width=\columnwidth]{fig_metrics_bar.pdf}
\caption{Cross-objective comparison at $\rho{=}1.5$. AMOS raises goodput and
deadline-hit rate while holding fidelity and utilization comparable---trading
surplus fidelity for schedule feasibility.}
\label{fig:bars}
\end{figure}

\subsection{Multi-objective operating point: weight sensitivity}
\label{sec:sens}
The knob $\lambda$ in~\eqref{eq:phi} moves the fidelity objective from pure
satisficing ($\lambda{=}0$) to pure maximization ($\lambda{=}1$).
Figure~\ref{fig:sens} and Table~\ref{tab:sens} sweep it. The result is a
\emph{robustness} finding rather than a wide Pareto front: across the entire
range, mean delivered fidelity rises by at most $+0.006$ and goodput changes by at
most $-0.005$ (both within seed noise). In this regime delivered fidelity of
successful tasks is pinned by device physics and the binding deadline constraint,
so weighting fidelity more heavily buys almost no delivered fidelity while
slightly reducing goodput at high load. This justifies the satisficing default
($\lambda{=}0$): it is at or near the best operating point, and the orchestrator's
behavior is not brittle to this weight. Where routing has more redundancy the knob
would gain leverage; we note this as a scope condition
(Section~\ref{sec:limitations}).

\begin{figure}[t]
\centering
\includegraphics[width=\columnwidth]{fig_sensitivity.pdf}
\caption{Operating-point sensitivity to the fidelity-weighting knob $\lambda$
(satisfice $\to$ maximize; mean $\pm$ s.d., 15 seeds). Both goodput (left axis)
and mean delivered fidelity (right axis) are essentially flat, so the operating
point is robust and the satisficing default is well-justified.}
\label{fig:sens}
\end{figure}

\begin{table}[t]
\caption{Weight-sensitivity of AMOS to the fidelity knob $\lambda$ (mean,
15 seeds). Variation is within seed noise.}
\label{tab:sens}
\centering
\small
\begin{tabular}{@{}l ccc ccc@{}}
\toprule
& \multicolumn{3}{c}{$\rho=1.0$} & \multicolumn{3}{c}{$\rho=1.5$}\\
\cmidrule(lr){2-4}\cmidrule(lr){5-7}
$\lambda$ & G & F & DL & G & F & DL\\
\midrule
0.0 & 0.518 & 0.468 & 0.996 & 0.489 & 0.473 & 0.995\\
0.5 & 0.515 & 0.470 & 0.996 & 0.485 & 0.471 & 0.996\\
1.0 & 0.513 & 0.474 & 0.996 & 0.490 & 0.474 & 0.996\\
\bottomrule
\end{tabular}
\end{table}

\subsection{Ablation: what does online learning add?}
Comparing AMOS and AMOS-S in Table~\ref{tab:metrics}, the two are statistically
indistinguishable across the stationary sweep ($0.518$ vs.\ $0.522$ at
$\rho{=}1.0$; $0.489$ vs.\ $0.491$ at $\rho{=}1.5$). We state this plainly:
\emph{when devices are stationary, the datasheet fidelity model already captures
the signal and online reliability learning adds no meaningful benefit over the
well-designed static policy.} The structured multi-objective placement, not the
learner, drives the advantage over baselines.

\subsection{Non-stationarity: when adaptation helps}
The learner's value appears when the datasheet goes stale. We repeat the
$\rho{=}1.0$ experiment but, after $40\%$ of the stream resolves, inflate the
realized error rate of the two devices datasheet-greedy policies most prefer by
$10\times$ (a persistent degradation event; the datasheet is \emph{not} updated).
Figure~\ref{fig:nonstat} plots windowed goodput and Table~\ref{tab:nonstat}
reports pre/post means. All policies drop at the change point, but AMOS recovers
the most: post-change goodput $0.406$ versus $0.395$ (AMOS-S) and $0.372$
(Max-Fidelity), i.e.\ the online derate cuts the post-degradation loss from
$0.152$ and $0.136$ to $0.122$. The effect is consistent but modest---a few
percentage points---which is the honest magnitude a lightweight reliability
estimator buys here.

\begin{figure}[t]
\centering
\includegraphics[width=\columnwidth]{fig_nonstat.pdf}
\caption{Windowed goodput under a mid-stream device-degradation event (dotted
line). Datasheet-greedy Max-Fidelity keeps routing to now-degraded devices; AMOS
learns to derate them and recovers a larger share of goodput.}
\label{fig:nonstat}
\end{figure}

\begin{table}[t]
\caption{Goodput before/after the degradation event ($\rho{=}1.0$, 15 seeds).}
\label{tab:nonstat}
\centering
\small
\begin{tabular}{@{}lccc@{}}
\toprule
Policy & Pre & Post & Drop\\
\midrule
Max-Fidelity & 0.524 & 0.372 & $-0.152$\\
AMOS-S       & 0.531 & 0.395 & $-0.136$\\
AMOS         & 0.527 & \textbf{0.406} & $\mathbf{-0.122}$\\
\bottomrule
\end{tabular}
\end{table}

\subsection{Summary of findings}
(i)~A deadline- and fidelity-aware satisficing orchestrator substantially
outperforms fidelity- and latency-greedy scheduling under contention, widening to
$\sim\!2\times$ at overload and near-parity when uncontended.
(ii)~Least-laxity ordering is the dominant contributor, sustaining $>\!99\%$
deadline-hit rate. (iii)~The operating point is robust to the fidelity-weighting
knob. (iv)~Online adaptation is negligible under stationarity but recovers a
consistent fraction of goodput after device degradation.

% ===================================================================
\section{Limitations and Threats to Validity}
\label{sec:limitations}
% ===================================================================
We are explicit about what this study does and does not establish.
\emph{First}, results are from simulation, not hardware; the fidelity model
\eqref{eq:mono}--\eqref{eq:dist} is a first-order abstraction omitting crosstalk,
correlated/non-Markovian noise, native-gate compilation, connectivity/SWAP
overhead, readout error, and realistic entanglement-distribution scheduling and
buffering. Absolute numbers are relative comparisons within a common model, not
device predictions. \emph{Second}, the workload regime was chosen so placement is
consequential (a $\approx\!0.70$ feasibility ceiling); where nearly all tasks are
trivially feasible or infeasible, orchestration matters less and gaps shrink.
\emph{Third}, the flat weight-sensitivity (Section~\ref{sec:sens}) is specific to
this regime: when most tasks have a single feasible host, the fidelity knob has
little routing leverage; richer device redundancy would widen the trade-off.
\emph{Fourth}, the learner is a tabular per-state/per-device estimator; richer
function approximation might extract more from non-stationarity at the cost of the
interpretability we relied on. \emph{Fifth}, we assume gang-scheduled distributed
execution on simultaneously free nodes and do not model an entanglement-generation
control plane; a joint circuit-cutting-plus-orchestration formulation is future
work. These caveats bound the claims to: \emph{within a transparent, reproducible
model, structured multi-objective orchestration is markedly more robust than
greedy baselines, its operating point is stable, and cheap online adaptation adds
a small, consistent benefit under drift.}

% ===================================================================
\section{Conclusion}
\label{sec:conclusion}
% ===================================================================
We framed noise-aware DQC as an AI-driven, multi-objective orchestration problem
and introduced AMOS, which orders by laxity, satisfices fidelity while optimizing
deadline feasibility, speed, and load, and adapts to calibration drift via an
online reliability derate. In a discrete-event simulator of a heterogeneous NISQ
fleet with entanglement-mediated distribution, AMOS matched the best baseline when
uncontended and improved goodput by up to $\sim\!2\times$ under overload, driven
mainly by least-laxity ordering and satisficing placement; its operating point was
robust to the fidelity-weighting knob, and online adaptation helped modestly and
only under non-stationarity. Future work includes hardware-calibrated noise
models, regimes with greater routing redundancy that expose a wider
multi-objective front, a joint circuit-cutting/orchestration formulation, an
explicit entanglement-distribution control plane, and richer learned reliability
models evaluated against the interpretable baseline established here.

% ===================================================================
\begin{thebibliography}{00}
\bibitem{caleffi2024survey} M.~Caleffi, M.~Amoretti, D.~Ferrari, J.~Illiano,
A.~Manzalini, and A.~S.~Cacciapuoti, ``Distributed quantum computing: A survey,''
\emph{Computer Networks}, vol.~254, p.~110672, 2024.
\bibitem{barral2025review} D.~Barral \emph{et al.}, ``Review of distributed
quantum computing: From single QPU to high performance quantum computing,''
\emph{Computer Science Review}, vol.~57, p.~100747, 2025.
\bibitem{chehimi2022qin} M.~Chehimi and W.~Saad, ``Physics-informed quantum
communication networks: A vision toward the quantum internet,'' \emph{IEEE
Network}, vol.~36, no.~5, pp.~32--38, 2022.
\bibitem{piveteau2024knitting} C.~Piveteau and D.~Sutter, ``Circuit knitting with
classical communication,'' \emph{IEEE Trans. Inf. Theory}, vol.~70, no.~4,
pp.~2734--2745, 2024.
\bibitem{du2024cutting} Z.~Du, W.~Zhang, W.~Wei, J.~Chen, T.~Han, Z.~Liang, and
Y.~Mao, ``Efficient circuit cutting and scheduling in a multi-node quantum system
with dynamic EPR pairs,'' \emph{arXiv preprint} arXiv:2412.18709, 2024.
\bibitem{carrera2024realtime} A.~Carrera~Vazquez, C.~Tornow, D.~Rist\`e,
S.~Woerner, M.~Takita, and D.~J.~Egger, ``Combining quantum processors with
real-time classical communication,'' \emph{Nature}, vol.~636, p.~75, 2024.
\bibitem{bahrani2024resource} S.~Bahrani, R.~D.~Oliveira, J.~M.~Parra-Ullauri,
R.~Wang, and D.~Simeonidou, ``Resource management and circuit scheduling for
distributed quantum computing interconnect networks,'' \emph{arXiv preprint}
arXiv:2409.12675, 2024.
\bibitem{nguyen2024drlq} H.~T.~Nguyen, M.~Usman, and R.~Buyya, ``DRLQ: A deep
reinforcement learning-based task placement for quantum cloud computing,'' in
\emph{Proc. 2024 IEEE 17th Int. Conf. on Cloud Computing (CLOUD)}, 2024,
pp.~475--481.
\bibitem{nguyen2025qfor} H.~T.~Nguyen, M.~Usman, and R.~Buyya, ``QFOR: A
fidelity-aware orchestrator for quantum computing environments using deep
reinforcement learning,'' \emph{ACM Trans. Quantum Computing}, 2025.
\bibitem{luo2025rl} W.~Luo, J.~Zhao, T.~Zhan, and Q.~Guan, ``Adaptive job
scheduling in quantum clouds using reinforcement learning,'' \emph{arXiv
preprint} arXiv:2506.10889, 2025.
\bibitem{chandra2025netops} N.~K.~Chandra \emph{et al.}, ``Network operations
scheduling for distributed quantum computing,'' \emph{arXiv preprint}
arXiv:2511.13687, 2025.
\bibitem{main2025optical} D.~Main, P.~Drmota, D.~P.~Nadlinger, E.~M.~Ainley,
A.~Agrawal, B.~C.~Nichol, R.~Srinivas, G.~Araneda, and D.~M.~Lucas,
``Distributed quantum computing across an optical network link,'' \emph{Nature},
vol.~638, no.~8050, pp.~383--388, 2025.
\bibitem{larasati2025ftdqc} H.~T.~Larasati and B.-S.~Choi, ``Towards
fault-tolerant distributed quantum computation (FT-DQC): Taxonomy, recent
progress, and challenges,'' \emph{ICT Express}, vol.~11, no.~3, pp.~417--435,
2025.
\bibitem{knorzer2025review} J.~Kn\"orzer, X.~Liu, B.~F.~Schiffer, and J.~Tura,
``Distributed quantum information processing: A review of recent progress,''
\emph{arXiv preprint} arXiv:2510.15630, 2025.
\bibitem{meddeb2025qinternet} A.~Meddeb, ``Quantum internet building blocks:
State of research and development,'' \emph{Computer Networks}, vol.~261,
p.~111151, 2025.
\bibitem{harada2024wire} H.~Harada, K.~Wada, and N.~Yamamoto, ``Doubly optimal
parallel wire cutting without ancilla qubits,'' \emph{PRX Quantum}, vol.~5,
no.~4, p.~040308, 2024.
\bibitem{schmitt2025cutting} L.~Schmitt, C.~Piveteau, and D.~Sutter, ``Cutting
circuits with multiple two-qubit unitaries,'' \emph{Quantum}, vol.~9, p.~1634,
2025.
\bibitem{harrow2025optimal} A.~W.~Harrow and A.~Lowe, ``Optimal quantum circuit
cuts with application to clustered Hamiltonian simulation,'' \emph{PRX Quantum},
vol.~6, p.~010316, 2025.
\bibitem{ferrari2023modular} D.~Ferrari, S.~Carretta, and M.~Amoretti, ``A
modular quantum compilation framework for distributed quantum computing,''
\emph{IEEE Trans. Quantum Eng.}, vol.~4, pp.~1--13, 2023.
\bibitem{escofet2023hungarian} P.~Escofet, A.~Ovide, C.~G.~Almud\'ever,
E.~Alarc\'on, and S.~Abadal, ``Hungarian qubit assignment for optimized mapping
of quantum circuits on multi-core architectures,'' \emph{IEEE Comput. Archit.
Lett.}, vol.~22, no.~2, pp.~161--164, 2023.
\bibitem{pastor2024partitioning} A.~Pastor, P.~Escofet, S.~B.~Rached,
E.~Alarc\'on, P.~Barlet-Ros, and S.~Abadal, ``Circuit partitioning for multi-core
quantum architectures with deep reinforcement learning,'' \emph{arXiv preprint}
arXiv:2401.17976, 2024.
\bibitem{ferrari2024execution} D.~Ferrari, M.~Bandini, and M.~Amoretti,
``Execution management of distributed quantum computing jobs,'' in \emph{Proc.
IEEE Int. Conf. Quantum Computing and Engineering (QCE)}, 2024, pp.~150--154.
\bibitem{nguyen2024review} H.~T.~Nguyen, P.~Krishnan, D.~Krishnaswamy, M.~Usman,
and R.~Buyya, ``Quantum cloud computing: A review, open problems, and future
directions,'' \emph{arXiv preprint} arXiv:2404.11420, 2024.
\bibitem{nguyen2024iquantum} H.~T.~Nguyen, M.~Usman, and R.~Buyya, ``iQuantum: A
toolkit for modeling and simulation of quantum computing environments,''
\emph{Softw. Pract. Exper.}, vol.~54, no.~6, pp.~1141--1171, 2024.
\end{thebibliography}

\end{document}
